% ============================================================
%  Список литературы
%  Оформление по ГОСТ Р 7.0.100-2018.
%  ВНИМАНИЕ: перед сдачей в издательство сверьте выходные данные
%  и ISBN с теми изданиями, которые реально используются в курсе, —
%  у переизданий они различаются.
% ============================================================
\chapter*{Список литературы}
\addcontentsline{toc}{chapter}{Список литературы}
\markboth{Список литературы}{}

\section*{Основная литература}

\begin{enumerate}[leftmargin=*,labelsep=6pt]

\item Хорстманн, К.\,С. Java. Библиотека профессионала. Том~1. Основы~/
К.\,С.~Хорстманн.~--- 12-е изд.~--- Москва: Диалектика, 2022.~--- 864~с.~---
Текст: непосредственный.

\item Хорстманн, К.\,С. Java. Библиотека профессионала. Том~2. Расширенные
средства программирования~/ К.\,С.~Хорстманн.~--- 12-е изд.~--- Москва:
Диалектика, 2022.~--- 176~с.~--- Текст: непосредственный.

\item Блох, Дж. Java. Эффективное программирование~/ Дж.~Блох.~--- 3-е изд.~---
Москва: Диалектика, 2019.~--- 464~с.~--- Текст: непосредственный.

\item Шилдт, Г. Java. Полное руководство~/ Г.~Шилдт.~--- 12-е изд.~--- Москва:
Диалектика, 2022.~--- 1344~с.~--- Текст: непосредственный.

\item Эккель, Б. Философия Java~/ Б.~Эккель.~--- 4-е изд.~--- Санкт-Петербург:
Питер, 2021.~--- 1168~с.~--- Текст: непосредственный.

\end{enumerate}

\section*{Дополнительная литература}

\begin{enumerate}[leftmargin=*,labelsep=6pt,resume]

\item Приёмы объектно-ориентированного проектирования. Паттерны проектирования~/
Э.~Гамма, Р.~Хелм, Р.~Джонсон, Дж.~Влиссидес.~--- Санкт-Петербург: Питер,
2021.~--- 448~с.~--- Текст: непосредственный.

\item Мартин, Р.\,С. Чистый код: создание, анализ и рефакторинг~/
Р.\,С.~Мартин.~--- Санкт-Петербург: Питер, 2019.~--- 464~с.~---
Текст: непосредственный.

\item Мартин, Р.\,С. Чистая архитектура. Искусство разработки программного
обеспечения~/ Р.\,С.~Мартин.~--- Санкт-Петербург: Питер, 2018.~--- 352~с.~---
Текст: непосредственный.

\item Фаулер, М. Рефакторинг. Улучшение проекта существующего кода~/
М.~Фаулер.~--- 2-е изд.~--- Москва: Диалектика, 2019.~--- 448~с.~---
Текст: непосредственный.

\item Уоллс, К. Spring в действии~/ К.~Уоллс.~--- 6-е изд.~--- Москва: ДМК Пресс,
2022.~--- 520~с.~--- Текст: непосредственный.

\item Фримен, Э. Head First. Паттерны проектирования~/ Э.~Фримен,
Э.~Робсон.~--- Санкт-Петербург: Питер, 2022.~--- 656~с.~---
Текст: непосредственный.

\item Гётц, Б. Java Concurrency на практике~/ Б.~Гётц, Т.~Пайерлс,
Дж.~Блох.~--- Санкт-Петербург: Питер, 2020.~--- 464~с.~---
Текст: непосредственный.

\item Чакон, С. Git для профессионального программиста~/ С.~Чакон,
Б.~Штрауб.~--- Санкт-Петербург: Питер, 2021.~--- 496~с.~---
Текст: непосредственный.

\end{enumerate}

\section*{Нормативные документы}

\begin{enumerate}[leftmargin=*,labelsep=6pt,resume]

\item ГОСТ~7.60--2003. Издания. Основные виды. Термины и определения:
межгосударственный стандарт: издание официальное: введён в действие
Постановлением Госстандарта России от 25.11.2003 №~331-ст.~--- Москва:
Стандартинформ, 2004.~--- Текст: непосредственный.

\item ГОСТ~Р~7.0.100--2018. Библиографическая запись. Библиографическое
описание. Общие требования и правила составления: национальный стандарт
Российской Федерации: издание официальное: утверждён и введён в действие
Приказом Росстандарта от 03.12.2018 №~1050-ст.~--- Москва: Стандартинформ,
2018.~--- Текст: непосредственный.

\end{enumerate}

\section*{Электронные ресурсы}

\begin{enumerate}[leftmargin=*,labelsep=6pt,resume]

\item The Java\textsuperscript{\textregistered} Language Specification:
Java SE~21 Edition: сайт.~--- URL: \url{https://docs.oracle.com/javase/specs/}
(дата обращения: 01.09.2026).~--- Текст: электронный.

\item The Java\textsuperscript{\textregistered} Virtual Machine Specification:
Java SE~21 Edition: сайт.~--- URL: \url{https://docs.oracle.com/javase/specs/}
(дата обращения: 01.09.2026).~--- Текст: электронный.

\item Java Platform, Standard Edition Documentation: сайт.~---
URL: \url{https://docs.oracle.com/en/java/javase/21/} (дата обращения:
01.09.2026).~--- Текст: электронный.

\item Spring Framework Documentation: сайт.~---
URL: \url{https://docs.spring.io/spring-framework/reference/}
(дата обращения: 01.09.2026).~--- Текст: электронный.

\item Spring Boot Reference Documentation: сайт.~---
URL: \url{https://docs.spring.io/spring-boot/} (дата обращения:
01.09.2026).~--- Текст: электронный.

\item JavaFX Documentation: сайт.~--- URL: \url{https://openjfx.io/}
(дата обращения: 01.09.2026).~--- Текст: электронный.

\item JUnit~5 User Guide: сайт.~---
URL: \url{https://junit.org/junit5/docs/current/user-guide/}
(дата обращения: 01.09.2026).~--- Текст: электронный.

\item Hibernate ORM Documentation: сайт.~---
URL: \url{https://hibernate.org/orm/documentation/} (дата обращения:
01.09.2026).~--- Текст: электронный.

\item Apache Maven Project: сайт.~--- URL: \url{https://maven.apache.org/}
(дата обращения: 01.09.2026).~--- Текст: электронный.

\item Pro Git~/ S.~Chacon, B.~Straub: сайт.~---
URL: \url{https://git-scm.com/book/ru/v2} (дата обращения: 01.09.2026).~---
Текст: электронный.

\end{enumerate}
